% ch3-part2.tex — Chapter 3, Sprints 4 & 5 (Section 3.4, Section 3.5).
% Continuation of Chapter 3; the \chapter command lives in ch3-part1.tex.
% Compile with XeLaTeX/LuaLaTeX + fontspec (unicode kept verbatim).

\clearpage
\section{Sprint 4: Productivity Suite}

\subsection{Sprint goal}

The first three sprints built the shared backbone: identity, projects, and the agile/task
engine. Sprint~4 turns to the surface each person touches every day on their own: a private
to-do list, the check-in/out attendance ritual, the shared calendar, and the reminders that
tie it all together. Four modules land in one sprint because they share a
\textbf{cron-driven, multi-channel notification backbone}: each runs a scheduled job that fans
the same four delivery channels (in-app push, email, Telegram, ntfy) out to a user's enabled
endpoints, the per-minute jobs guarded by a single PostgreSQL-backed distributed lock. Where Sprint~3's
authorization was project-scoped and role-graduated, most of this sprint is personal: the
to-do, calendar and attendance data is \texttt{.own}-scoped and enforced in the data layer. At
\textbf{44 story points} it is a broad rather than deep sprint: four self-contained aggregates,
each modest on its own, unified by the notification backbone.

\subsection{Sprint backlog}

The user stories for this sprint, repeated from the product backlog in Section~2.2 so the
sprint reads self-contained:

\begin{longtable}{@{}l>{\raggedright\arraybackslash}p{0.50\textwidth}lll@{}}
\caption{Sprint 4 backlog.}\label{tab:s4-backlog-detail} \\
\tableheaderrow
\thcell{US-ID} & \thcell{Story} & \thcell{SP} & \thcell{Est.} & \thcell{Priority} \\
\midrule
\endfirsthead
\caption[]{Sprint 4 backlog (continued).} \\
\tableheaderrow
\thcell{US-ID} & \thcell{Story} & \thcell{SP} & \thcell{Est.} & \thcell{Priority} \\
\midrule
\endhead
\bottomrule
\endfoot
\rowcolor{rowalt}
US-S4-01 & As a user, I want a private to-do list with sub-tasks, priorities, and statuses, so that I track my own work separately from project tasks. & 5 & 2.5d & Should \\
US-S4-02 & As a user, I want due/reminder dates on my to-dos that notify me across my channels, so that nothing slips. & 5 & 2d & Should \\
\rowcolor{rowalt}
US-S4-03 & As an employee, I want to check in and out (remote/onsite) to open and close work sessions, so that attendance is recorded. & 8 & 4d & Must \\
US-S4-04 & As an employee, I want my worked time computed per business day, so that I can see my hours. & 3 & 1.5d & Should \\
\rowcolor{rowalt}
US-S4-05 & As a manager, I want per-user and per-team attendance statistics, so that I can review presence. & 5 & 2.5d & Should \\
US-S4-06 & As a user, I want to create calendar events and meetings with participants, so that schedules are shared. & 8 & 3.5d & Should \\
\rowcolor{rowalt}
US-S4-07 & As a user, I want multi-channel reminders before each event, so that I attend on time. & 3 & 1.5d & Should \\
US-S4-08 & As a user, I want to schedule project and personal reminders delivered on my channels, so that I am nudged at the right moment. & 5 & 2.5d & Should \\
\rowcolor{rowalt}
US-S4-09 & As a user, I want to see and dismiss my pending reminders, so that I can manage them. & 2 & 1d & Could \\
\midrule
 & \textbf{Sprint 4 subtotal} & \textbf{44} & \textbf{4w} & \\
\end{longtable}

\subsection{Requirements analysis}

The sprint delivers four modules. The features, actors, and use cases are collected in
Figure~\ref{fig:3-23}. Two actor archetypes appear: a \textbf{User} (any authenticated employee)
who manages their own to-dos, attendance, calendar events, and reminders, and a
\textbf{Manager} (CEO, CTO, CMO, or Project Manager) who additionally reviews their team's
attendance and schedules project-scoped reminders. The shared ``notify across channels'' use
case is drawn explicitly because it is the seam that binds the sprint: all four
modules \texttt{<<include>>} it from their reminder paths.

% Figure 3.23 — render diagrams-src/fig_3_23.puml -> figures/fig_3_23.(png|pdf)
{\renewcommand{\figmaxheight}{0.75\textheight}%
\reportfigure[\textwidth]{figures/fig_3_23}{Sprint 4 use-case diagram.}{3-23}}

\subsubsection{Module A --- Personal To-Dos}

\textbf{Functional description.} This module is a private per-user checklist, kept deliberately
separate from the project \texttt{Task} engine of Sprint~3 --- its own tables, its own status and
priority enums (\texttt{Pending}/\texttt{InProgress}/\texttt{Completed},
\texttt{Low}/\texttt{Medium}/\texttt{High}), so personal and project tasks can evolve independently. A
to-do can carry sub-tasks through a self-relation (cascade-deleted with the parent), favourites,
an archive flag, a due date, an optional reminder date, manual drag-ordering via
\texttt{displayOrder}, attachments and comments. The module is feature-rich without new
infrastructure because it reuses the existing upload, locking, and notification modules. Every
read and write is scoped by \texttt{userId} in the repository layer --- the \texttt{.own} permissions
are enforced in the data query, not by the guard --- so a user cannot fetch another user's list.

\textbf{Actors.} A \textbf{User} (all authenticated roles) creates, reads, updates, deletes, and
reorders their own personal tasks and comments. The module has no manager surface; every
personal-task permission carries the \texttt{.own} suffix and is granted to all roles through
\texttt{DEFAULT\_PERMISSIONS\_FOR\_ALL\_ROLES}.

\textbf{Business rules.}
\begin{itemize}
  \item Every read and write filters by \texttt{userId} in the repository, enforcing
    \textbf{strict read isolation}: \texttt{findUniqueOrThrow(\{ id, userId \})} returns 404 for
    any task the caller does not own.
  \item Sub-tasks are a \textbf{self-referencing relation} on \texttt{parentTaskId} with
    \texttt{onDelete: Cascade}, so deleting a parent removes its children atomically.
  \item Multilingual content follows the platform-wide i18n content-table pattern
    (\texttt{UserTaskContent} keyed by \texttt{@@unique([taskId, language])}).
  \item If file uploads fail after the database write, the service rolls back by deleting the
    already-uploaded files --- a cleanup guard not present in the project task module.
  \item The \texttt{reminderDate} field is optional; once the per-minute cron detects it has
    passed, it fans out to the owner's enabled channels and marks \texttt{notified = true} so the
    reminder fires exactly once.
\end{itemize}

\textbf{Reminder cron sequence.} The one background behaviour worth a diagram is the due-date
engine (Figure~\ref{fig:3-24}). A per-minute \texttt{@Cron} acquires a Postgres-backed distributed
lock (55-second TTL) so only one instance runs it, fetches to-dos whose \texttt{reminderDate} has
passed and are still un-notified, fans each one out to the owner's enabled channels, and marks
\texttt{notified = true}. It is the same lock-then-scan shape all four Sprint~4 crons use.

% Figure 3.24 — render diagrams-src/fig_3_24.puml -> figures/fig_3_24.(png|pdf)
\reportsideways{figures/fig_3_24}{Personal-task reminder cron sequence.}{3-24}

\textbf{Class diagram.} Figure~\ref{fig:3-25} shows a small, cohesive aggregate: a
\texttt{UserTask} root with the content-table split for i18n, comments carrying their own author
\texttt{userId}, and attachments stored as disk paths.

% Figure 3.25 — render diagrams-src/fig_3_25.puml -> figures/fig_3_25.(png|pdf)
\reportfigure{figures/fig_3_25}{Personal To-Dos class diagram.}{3-25}

\subsubsection{Module B --- Time \& Attendance}

\textbf{Functional description.} Attendance is enforced as a daily ritual. On entering any
authenticated dashboard route the user hits a full-screen gate that blocks the app until they
check in --- choosing Remote or Onsite --- or explicitly ``join as viewer''. Checking in opens a
\texttt{WorkSession}; checking out closes the latest open session, computes worked minutes, and
optionally captures a mood and a journey note. Sessions roll up under one \texttt{WorkDay} per
user per business day, which also holds day-level attributes (mood, a manager's performance
rating, notes). Beyond recording presence the module nudges for compliance --- late starts and
forgotten sessions trigger notifications --- and reports per-user and per-manager attendance
statistics.

\textbf{Actors.} A \textbf{Worker} (all authenticated roles) checks in and out, views their own
work days and statistics, and adds daily mood or notes. A \textbf{Manager} (CEO, CTO, CMO,
Project Manager) reviews team attendance, views per-user and per-team statistics, and adds a
performance rating or manager notes --- access gated by
\texttt{canUserManageUsers}/\texttt{canUserManageTeams} inside the service.

\textbf{Business rules.}
\begin{itemize}
  \item \textbf{Only one open session per day}: the service rejects a check-in if the current
    \texttt{WorkDay} already contains a session with no \texttt{endTime}.
  \item \textbf{Late-start detection} fires at two thresholds --- morning (08:45~UTC) and
    afternoon (12:30~UTC) --- with a single-notification-per-half-day guard to prevent duplicate
    alerts.
  \item A nightly \textbf{auto-close cron} at 03:00~UTC closes all open sessions, capping each
    session's \texttt{endTime} at shift end (16:15~UTC weekdays, 12:30~UTC Saturdays) rather than
    at 03:00, so workers are not penalized for after-hours time.
  \item \textbf{Workers cannot modify} \texttt{endTime} or \texttt{timeSpentInMinutes} on their
    own sessions; the service strips those fields before persisting.
  \item Manager statistics endpoints correctly add a \texttt{canUserManageUsers} check before
    querying, so a manager sees only the users they actually manage (business-unit-scoped).
  \item Location and device are recorded per session via \texttt{WorkSessionLocation}
    (\texttt{REMOTE}/\texttt{ONSITE}) and \texttt{WorkSessionDevice}
    (\texttt{DESKTOP}/\texttt{MOBILE}/\texttt{TABLET}/\texttt{OTHER}).
\end{itemize}

\textbf{Check-in sequence.} Figure~\ref{fig:3-26} shows the check-in path: the service resolves
the current business day (creating one if none exists), verifies no session is already open,
opens a new \texttt{WorkSession} with the chosen location and device, and returns the updated
work day.

% Figure 3.26 — render diagrams-src/fig_3_26.puml -> figures/fig_3_26.(png|pdf)
\reportsideways{figures/fig_3_26}{Check-in sequence diagram.}{3-26}

\textbf{Class diagram.} Figure~\ref{fig:3-27} shows the two-level \texttt{WorkDay} \textrightarrow{}
\texttt{WorkSession} split. The single documented seam into Sprint~3 is visible: a
\texttt{WorkSession} owns the \texttt{TaskTimeEntry} rows the Tasks module attaches to it, so
logged task time is anchored to an attendance session.

% Figure 3.27 — render diagrams-src/fig_3_27.puml -> figures/fig_3_27.(png|pdf)
{\renewcommand{\figmaxheight}{0.60\textheight}%
\reportfigure[0.90\textwidth]{figures/fig_3_27}{Time \& Attendance class diagram.}{3-27}}

\subsubsection{Module C --- Events \& Calendar}

\textbf{Functional description.} This is the shared company calendar. Three kinds of entry are
distinguished by an \texttt{EventType} enum --- \texttt{Meeting} and \texttt{Event}
(organizational, targetable at named participants or the whole company via \texttt{toAllUsers})
and \texttt{PersonalEvent} (a private per-user entry) --- each surfaced on its own calendar page.
The calendar is a custom month/week/day/agenda view built on \texttt{date-fns} and
\texttt{@dnd-kit} drag-and-drop. The module's second concern is reminders: a per-minute cron
computes \emph{escalating} notification thresholds --- 24~hours out, then 2.5~hours, then
30~minutes, then a final 15-minute window --- and fans each one out to the resolved audience
across the four delivery channels.

\textbf{Actors.} A \textbf{User} (all authenticated roles) creates, reads, and manages their own
events (including personal events). An \textbf{Executive} (CEO, CTO, CMO) can additionally
modify or delete organizational events created by others. Read visibility is enforced in the
query: a user sees only events where \texttt{toAllUsers} is true, they created it, or they are a
participant.

\textbf{Business rules.}
\begin{itemize}
  \item \textbf{Escalating notification thresholds}: the cron evaluates remaining time to the
    event start and advances \texttt{nextNotificationTime} through four bands (24h, 2.5h, 30min,
    15min). Once remaining time drops below 25~minutes, \texttt{isNotified} is set to \texttt{true}
    and no further notifications fire.
  \item The \texttt{EventParticipant} join uses a \textbf{composite primary key}
    \texttt{@@id([eventId, userId])}, preventing duplicate invitations.
  \item A \texttt{PersonalEvent} is editable and deletable \textbf{only by its creator}; the
    service checks \texttt{event.createdById === req.user.id}.
  \item \texttt{EventContent} follows the platform-wide i18n content-table pattern, consistent
    with \texttt{ProjectContent} and \texttt{SprintContent} from earlier sprints.
  \item Event colours are constrained to a six-value enum (\texttt{Sky}, \texttt{Amber},
    \texttt{Violet}, \texttt{Rose}, \texttt{Emerald}, \texttt{Orange}).
\end{itemize}

\textbf{Event reminder cron.} Figure~\ref{fig:3-28} shows the escalating reminder cron: on each
tick it queries events whose \texttt{nextNotificationTime} has passed, resolves the audience
(creator, explicit participants, or all users if \texttt{toAllUsers}), fans out to each
recipient's enabled channels, computes the next threshold, and updates the event row.

% Figure 3.28 — render diagrams-src/fig_3_28.puml -> figures/fig_3_28.(png|pdf)
\reportsideways{figures/fig_3_28}{Event reminder cron with escalating thresholds.}{3-28}

\textbf{Class diagram.} Figure~\ref{fig:3-29} shows a clean three-table aggregate ---
\texttt{Event}, its i18n \texttt{EventContent}, and the \texttt{EventParticipant} join. The
denormalized \texttt{nextNotificationTime}/\texttt{isNotified} fields on \texttt{Event} let the
cron find events due to notify with a single indexed scan rather than recomputing per event
each minute.

% Figure 3.29 — render diagrams-src/fig_3_29.puml -> figures/fig_3_29.(png|pdf)
\reportfigure{figures/fig_3_29}{Events \& Calendar class diagram.}{3-29}

\subsubsection{Module D --- Reminders}

\textbf{Functional description.} Reminders is the sprint's connective tissue: a time-triggered,
multi-channel reminder table that other modules also write into. It serves two workflows.
\textbf{Manual} reminders let a project manager schedule a nudge for a specific user, linked to
a task, sprint, milestone, project, or free-form \texttt{CUSTOM} target, delivered over chosen
channels at a future time. \textbf{Automatic} reminders are created by the Tasks, Sprints and
Milestones services when those artefacts are created (the due-date nudges seen in Sprint~3),
plus crons that scan for overdue tasks (due date passed) and stuck tasks (no update for five
days). A per-minute scheduler is the delivery engine: it finds \texttt{PENDING} reminders whose
\texttt{reminderAt} has passed and dispatches them. Authorization mirrors Projects exactly ---
the same BusinessUnit-scoped executive model (CEO global, CTO TawerDev, CMO TawerCreative)
plus project membership --- which keeps the codebase consistent.

\textbf{Actors.} A \textbf{Team Member} (all project members) can read reminders for their
projects and dismiss their own reminders. A \textbf{Manager / Executive} (PM, PO, CEO, CTO, CMO)
can create, update, and delete reminders for projects they manage, following the same
business-unit-scoped access model introduced in Sprint~2.

\textbf{Business rules.}
\begin{itemize}
  \item \texttt{reminderAt} must be \textbf{in the future}; the service rejects any value
    \texttt{<= now()}.
  \item The reminder follows a \textbf{status machine}: \texttt{PENDING \textrightarrow{} SENT
    \textrightarrow{} DISMISSED} for normal flow, or \texttt{PENDING \textrightarrow{} CANCELLED}
    for manual cancellation. Dismissing an already-dismissed or cancelled reminder is rejected.
  \item The target is \textbf{polymorphic}: \texttt{entityType} (\texttt{TASK}, \texttt{SPRINT},
    \texttt{MILESTONE}, \texttt{PROJECT}, \texttt{CUSTOM}) plus \texttt{entityId} links a reminder
    to any entity kind without requiring a typed foreign key per entity.
  \item Two distinct \texttt{User} foreign keys separate \emph{who is reminded}
    (\texttt{userId}) from \emph{who scheduled it} (\texttt{createdById}).
  \item Overdue and stuck task detection is \textbf{idempotent}: the cron checks for an existing
    \texttt{PENDING}/\texttt{SENT} reminder with the same entity and message before creating a new
    one, preventing duplicate notifications.
  \item Default reminders are seeded automatically: one day before a task's due date, one day
    before sprint start and two days before sprint end, and seven days plus three days before a
    milestone's due date.
\end{itemize}

\textbf{Pending-reminder delivery sequence.} Figure~\ref{fig:3-30} shows the core delivery path:
the per-minute scheduler acquires the distributed lock, queries \texttt{PENDING} reminders whose
\texttt{reminderAt} has passed, resolves each recipient's channel preferences, fans out to the
enabled channels, and marks the reminder \texttt{SENT}.

% Figure 3.30 — render diagrams-src/fig_3_30.puml -> figures/fig_3_30.(png|pdf)
\reportsideways{figures/fig_3_30}{Pending-reminder delivery sequence.}{3-30}

\textbf{Class diagram.} Figure~\ref{fig:3-31} shows the two design decisions worth noting: the
dual-user foreign keys (\texttt{userId} vs \texttt{createdById}) and the polymorphic target
(\texttt{entityType} + \texttt{entityId}).

% Figure 3.31 — render diagrams-src/fig_3_31.puml -> figures/fig_3_31.(png|pdf)
\reportfigure{figures/fig_3_31}{Reminders class diagram.}{3-31}

\subsection{Implementation}

The sprint delivered four backend modules and the client screens that consume them.

The \textbf{personal to-do surface} comprises endpoints on
\texttt{PersonalTasksController} covering CRUD, sub-tasks, reorder, favourites,
archive, comments, and attachments --- all scoped by the caller's \texttt{userId}.

The \textbf{attendance surface} comprises endpoints on \texttt{WorkDaysController}:
check-in, check-out, current work day, own and manager work-day lists, worker and manager
updates, and four statistics endpoints (overview/details, each for own and manager views).

The \textbf{calendar surface} comprises endpoints on \texttt{EventsController}: event CRUD,
participant management, and the three event-type list queries (meetings, events, personal
events). The calendar is a bespoke month/week/day/agenda component built on \texttt{date-fns}
and \texttt{@dnd-kit}.

The \textbf{reminders surface} comprises endpoints across two controllers:
\texttt{RemindersController} (project-scoped CRUD) and \texttt{UserRemindersController}
(personal \texttt{/reminders/me} list and dismiss action).

On the \textbf{client}, the to-do list lives under a two-tab page (Personal / Project),
attendance is a full-screen gate wrapping the dashboard plus a header check-out button, the
calendar is three pages sharing one custom calendar component, and the reminders UI is embedded
inside the project-detail view.

\reportscreenshot{figures/screenshot_P23}{Personal to-do list with filters and task-detail sheet.}{P23}

\reportscreenshot{figures/screenshot_P24}{Check-in gate.}{P24}

\reportscreenshot{figures/screenshot_P25}{Calendar month view.}{P25}

\reportscreenshot{figures/screenshot_P26}{Project-detail reminders list.}{P26}

\subsection{Validation and testing}

The testing strategy for this sprint focused on two areas: \textbf{property-based frontend
testing} and \textbf{acceptance scenario verification}.

Two property-based suites (\texttt{fast-check} + \texttt{vitest}) pin client-side invariants
across the sprint's transform and parsing layers: the reminder Zod schema invariants and the
caster round-trip. Property-based testing generates hundreds of random inputs per run and
asserts an invariant holds across all of them, which is a stronger guarantee than a handful of
example cases.

The acceptance scenarios below record the Given/When/Then conditions each user story had to
satisfy:

\begin{longtable}{@{}p{1.8cm}p{8.3cm}p{3.8cm}@{}}
\caption{Sprint 4 acceptance scenarios.}\label{tab:s4-acceptance} \\
\tableheaderrow
\thcell{US-ID} & \thcell{Acceptance scenario (Given / When / Then)} & \thcell{Result} \\
\midrule
\endfirsthead
\caption[]{Sprint 4 acceptance scenarios (continued).} \\
\tableheaderrow
\thcell{US-ID} & \thcell{Acceptance scenario (Given / When / Then)} & \thcell{Result} \\
\midrule
\endhead
\bottomrule
\endfoot
\rowcolor{rowalt}
US-S4-01 & Given a user, When they create a to-do with sub-tasks/priority/status, Then it is stored under their \texttt{userId}, isolated from other users and from project tasks, with sub-tasks cascade-linked. & \pass{} Verified \\
US-S4-02 & Given a to-do with a \texttt{reminderDate}, When the time passes, Then a locked per-minute cron notifies the owner across enabled channels and marks it notified once. & \pass{} Verified \\
\rowcolor{rowalt}
US-S4-03 & Given an employee, When they check in (remote/onsite) and later check out, Then a \texttt{WorkSession} opens and closes with worked minutes computed; only one session is open at a time. & \pass{} Verified \\
US-S4-04 & Given a closed session, When the day is viewed, Then worked time is computed per business day. & \pass{} Verified \\
\rowcolor{rowalt}
US-S4-05 & Given a manager, When they open attendance statistics for their users/teams, Then aggregates and per-day hours are returned for the managed users, gated by \texttt{canUserManageUsers}. & \pass{} Verified \\
US-S4-06 & Given a user, When they create an event/meeting with participants, Then it is stored with content and participants and appears on the shared calendar with correct read visibility (creator, participant, or \texttt{toAllUsers}). & \pass{} Verified \\
\rowcolor{rowalt}
US-S4-07 & Given an upcoming event, When each threshold is reached, Then a reminder fans out over enabled channels at 24h/2.5h/30min/15min, advancing \texttt{nextNotificationTime} until \texttt{isNotified} is set. & \pass{} Verified \\
US-S4-08 & Given a manager, When they schedule a project or personal reminder for a user with channels and a future time, Then it is stored \texttt{PENDING} and delivered by the per-minute scheduler across enabled channels. & \pass{} Verified \\
\rowcolor{rowalt}
US-S4-09 & Given a user with pending reminders, When they open their list and dismiss one, Then it is marked dismissed (ownership-enforced via \texttt{userId} scoping). & \pass{} Verified \\
\end{longtable}

\subsection{Sprint review}

\textbf{Objectives achieved.} The sprint met its goal: the platform has a private task manager
with per-user isolation, an enforced attendance ritual with compliance nudges and manager
reporting, a shared calendar with escalating multi-channel reminders, and a project-aware
reminder engine that other modules write into. All four modules sit on one shared
infrastructure contribution: a cron-driven, lock-guarded, multi-channel notification backbone
where each per-minute cron acquires a PostgreSQL distributed lock (\texttt{SELECT \ldots{} FOR
UPDATE SKIP LOCKED} with a 55-second TTL), scans for pending items, fans out to the user's
enabled channels, and marks the item processed. This backbone is exactly what Sprint~5 will formalize
into a first-class notification module.

\textbf{Remaining work.} Hardening items identified during the sprint --- centralizing the
duplicated four-channel fan-out into a single dispatcher service, backend tests for the cron
timing logic, and tightening several access-control scopes --- are logged and prioritized for
future work (Annex~A).

\subsection{Sprint retrospective}

\textbf{Difficulties encountered.} The sprint's infrastructure challenge was designing the
PostgreSQL-based distributed lock without a Redis dependency. The lock had to guarantee that
only one API instance processes a given cron job in a horizontally-scaled deployment, survive
process crashes via a 55-second TTL, and be reusable across all four modules. The escalating
event notification thresholds (24h / 2.5h / 30min / 15min) required careful time arithmetic,
and the attendance module's business-day boundaries --- weekday shift-end at 16:15~UTC vs
Saturday half-day at 12:30~UTC, morning and afternoon late thresholds --- added further
time-domain complexity. Coordinating four modules that each produce notifications but share one
delivery backbone demanded establishing a consistent pattern before any module could ship.

\textbf{Lessons learned.} The lock-then-scan cron pattern proved highly reusable: once
established in the to-do reminder cron, the same structure (acquire lock, query pending items,
fan out, mark processed) was adopted by the event, reminder, and attendance crons with minimal
adaptation. Gating each channel on the user's preference booleans independently --- rather than
a single ``notify me'' flag --- gave fine-grained control without complicating the sender logic.
The polymorphic reminder target (\texttt{entityType} + \texttt{entityId}) also proved its value:
a single table and delivery path serve tasks, sprints, milestones, projects, and custom
reminders without per-entity-type plumbing.

\textbf{Improvements planned.} The four-channel fan-out logic is duplicated across several
consumers; a single \texttt{NotificationDispatcherService} would centralize channel selection
and delivery. Backend tests for the cron timing logic (threshold bands, business-day
boundaries) are needed. These items are carried forward as prioritized future work (Annex~A).

\subsection{Cumulative class diagram}

Figure~\ref{fig:3-32} focuses on the entities introduced in this sprint, following the
convention of the earlier cumulative views: the anchor classes from previous sprints ---
\texttt{User} (Sprint~1), \texttt{Project} (Sprint~2), and \texttt{Task}, \texttt{Milestone} and
\texttt{TaskTimeEntry} (Sprint~3) --- are shown in plain style and compacted to their
identifying fields, so the attachment points are clear, while this sprint's four new
aggregates are highlighted and grouped by module. The
seams are what make the picture cohere: every new aggregate hangs off \texttt{User}; a
\texttt{WorkSession} owns the \texttt{TaskTimeEntry} rows from Sprint~3 (attendance
anchors logged task time); and a \texttt{Reminder} reaches back into \texttt{Project},
\texttt{Task} and \texttt{Milestone} through its \texttt{projectId} and its polymorphic
\texttt{entityId}, tying the new personal layer to the project engine built in Sprints~2 and~3.

% Figure 3.32 — render diagrams-src/fig_3_32.puml -> figures/fig_3_32.(png|pdf)
\reportsideways{figures/fig_3_32}{Cumulative class diagram after Sprint 4.}{3-32}

\subsection{Sprint conclusion}

Sprint~4 delivered the everyday personal surface --- private tasks, enforced attendance, a
shared calendar, and a project-aware reminder engine --- unified by a cron-driven, lock-guarded,
multi-channel notification backbone. The lock-then-scan pattern and per-user channel gating
introduced here become the delivery substrate that Sprint~5 formalizes into a first-class
notification module.

\clearpage
\section{Sprint 5: Communication \& Operations}

\subsection{Sprint goal}

Sprint 4 left every module writing to the same four delivery channels through copy-pasted
fan-out code. Sprint 5 makes that backbone a first-class subject in its own right and adds the one consumer
that exists only to feed it. Two modules land together because they are the two halves of one concern:
\emph{reaching people}. The \textbf{notifications} module owns the delivery substrate: the in-app inbox, FCM push, the
per-user channel settings, and the Telegram/ntfy account links that every other module reads before it sends.
The \textbf{infrastructure monitoring} module is the platform watching itself: it registers the company's servers
and the services running on them, health-checks them every minute, and turns any outage or looming
subscription-expiry into an alert routed to the responsible managers over those same four channels. Putting them in one sprint is deliberate: monitoring is the clearest end-to-end demonstration of the
notification backbone, because an alert travels from a cron detecting a dead ping all the way to a Telegram
message, reusing the exact settings the notifications module defines. At \textbf{35 story points} it is a focused
sprint: one module that is broad but shallow (many channels, thin logic each) and one that is operationally
deep (six crons, an outbox, escalating expiry bands).

\subsection{Sprint backlog}

The stories drawn from the product backlog (§2.2) for this sprint:

\begin{longtable}{@{}l>{\raggedright\arraybackslash}p{0.50\textwidth}lll@{}}
\caption{Sprint 5 backlog.}\label{tab:s5-backlog-detail} \\
\tableheaderrow
\thcell{US-ID} & \thcell{Story} & \thcell{SP} & \thcell{Est.} & \thcell{Priority} \\
\midrule
\endfirsthead
\caption[]{Sprint 5 backlog (continued).} \\
\tableheaderrow
\thcell{US-ID} & \thcell{Story} & \thcell{SP} & \thcell{Est.} & \thcell{Priority} \\
\midrule
\endhead
\bottomrule
\endfoot
US-S5-01 & As a user, I want an in-app notification inbox (bell + list), so that I see system events in one place. & 5 & 2.5d & Must \\
US-S5-02 & As a user, I want push notifications (FCM) on my devices, so that I am alerted while away from the page. & 5 & 2d & Should \\
US-S5-03 & As a user, I want to configure my delivery channels (email, push, Telegram, ntfy), so that I control how I am reached. & 3 & 1.5d & Should \\
US-S5-04 & As a user, I want to link my Telegram account, so that I can receive alerts there. & 3 & 1.5d & Could \\
US-S5-05 & As a DevOps engineer, I want to register servers and services, so that infrastructure is inventoried. & 5 & 2.5d & Should \\
US-S5-06 & As the platform, I want to health-check servers (ICMP) and services (HTTP) every minute, so that outages are detected automatically. & 8 & 3.5d & Must \\
US-S5-07 & As a manager, I want multi-channel alerts when a server or service goes down, so that I can react quickly. & 5 & 2d & Must \\
US-S5-08 & As an operator, I want a public \texttt{/health} endpoint, so that the API's liveness is externally checkable. & 1 & 0.5d & Could \\
\midrule
 & \textbf{Sprint 5 subtotal} & \textbf{35} & \textbf{3w} & \\
\end{longtable}

\subsection{Requirements analysis}

The sprint use-case diagram (Figure~\ref{fig:3-33}) shows the actor split. A plain \textbf{User} manages their own inbox and
channel preferences; a \textbf{DevOps engineer} and the \textbf{CTO} register and manage the infrastructure, with the CTO
holding global write and the DevOps engineer scoped to servers they manage. A separate \textbf{Scheduler}
actor stands for the automated detectors --- no human triggers a health check --- and its alert path
\texttt{<<include>>}s the same multi-channel notification the human-facing features use, which is the point of pairing
the two modules.

% Figure 3.33 — render diagrams-src/fig_3_33.puml -> figures/fig_3_33.(png|pdf)
\reportfigure{figures/fig_3_33}{Sprint 5 use-case diagram.}{3-33}

\subsubsection{Module A — Notifications}

\textbf{What the module does.} This is the delivery substrate the whole platform sits on, and its most important
design fact is a negative one: there is \textbf{no central dispatcher}. The \texttt{NotificationsModule} owns only two of
the four channels --- the in-app inbox (persisted rows) and FCM web push --- and it exposes one shared method,
\texttt{createNotificationFromSystem(userId, payload)}, that seven other modules call to reach a single user in-app and
by push. The remaining three channels --- email, Telegram, ntfy --- are \emph{not} orchestrated here.
Each consumer module reads the target user's \texttt{UserNotificationSettings}, \texttt{UserTelegramBot} and
\texttt{UserNtfyIntegration}, then calls \texttt{MailService}, \texttt{TelegramService} or \texttt{NtfyService} directly, gating each
channel on the matching boolean flag. The consequence is that the four-channel selection logic is
copy-pasted across roughly seven consumers rather than living in one place --- high cohesion inside the module,
but a cross-cutting fan-out that has been duplicated rather than centralized, which constitutes the subsystem's
main architectural weakness.

Two entry points exist. A human-driven \textbf{broadcast} --- \texttt{POST /notifications} --- lets an authenticated user send
a titled, optionally imaged notification to a list of users or to everyone, persisting it in-app and pushing it
by FCM. And the \textbf{system-driven} per-user path already described, which crons and services use for
the reminders and nudges of Sprints 3 and 4. All six controller routes sit behind \texttt{HasPermissionGuard}, and
reads/updates/deletes are correctly scoped to the caller (a user cannot read, mark, or delete another user's
notifications), which is a genuine strength. The channel defaults encode the intended UX: email and
push are opt-out (default on), Telegram and ntfy are opt-in (default off), because the latter two need extra
per-user setup. The component view of the fan-out is Figure~\ref{fig:3-34}; a single system notification
crossing the channels is Figure~\ref{fig:3-35}.

% Figure 3.34 — render diagrams-src/fig_3_34.puml -> figures/fig_3_34.(png|pdf)
\reportfigure{figures/fig_3_34}{Multi-channel delivery pipeline component view.}{3-34}

% Figure 3.35 — render diagrams-src/fig_3_35.puml -> figures/fig_3_35.(png|pdf)
\reportfigure{figures/fig_3_35}{Single system notification across all channels.}{3-35}

The class slice (Figure~\ref{fig:3-36}) is the module's seven tables. The shape worth reading is the split between the
notification \emph{envelope} (\texttt{Notification}), its translatable \emph{content} (\texttt{NotificationContent}, in practice
English-only --- the multi-language capability is dormant, consistent with the codebase-wide i18n finding), and
the per-recipient \emph{inbox row} (\texttt{UserNotification}), whose \texttt{@@unique([notificationId, userId])} makes fan-out
re-sends idempotent per user. The three settings/integration tables hang one-to-one off \texttt{User} and are
the configuration every consumer reads before it sends.

% Figure 3.36 — render diagrams-src/fig_3_36.puml -> figures/fig_3_36.(png|pdf)
\reportfigure{figures/fig_3_36}{Notifications class slice with channel-configuration tables.}{3-36}

\textbf{Limitations.} The module reads cleanly and is injection-safe (Prisma-only) and ownership-correct on
read/update/delete, but four verified issues qualify how ``done'' it is. \textbf{ntfy is effectively dead:} the
settings UI instructs the user to subscribe to a topic equal to their \texttt{user.id}, but no code path ever writes
that value into \texttt{UserNtfyIntegration.topic}, and \texttt{NtfyService} early-returns on a falsy topic --- so enabling
ntfy in the UI silently delivers nothing. \textbf{The in-app inbox is coupled to the push toggle:}
because the \texttt{UserNotification} rows are persisted inside \texttt{createNotificationFromSystem}, which consumers gate on
\texttt{pushNotificationsEnabled}, a user who turns push off also loses their in-app history for system notifications ---
the two are not independently controllable. \textbf{Broadcast is open to every role:} \texttt{notification.create}
sits in the default permission set granted to all 31 roles with no service-level restriction, so any
authenticated user can \texttt{sendToAllUsers} and spam the whole company in-app and by push, with no rate limit. And the \textbf{channel error contract is inconsistent} --- \texttt{MailService} re-throws on
failure while Firebase, Telegram and ntfy swallow, so a consumer that fans out to all four can get an unhandled
rejection only from mail. Testing is skeletal: the two notification specs are \texttt{toBeDefined()}
stubs and the service spec cannot even resolve its DI graph; only \texttt{mail.service.spec.ts} is a real test.

\subsubsection{Module B — Infrastructure Monitoring}

\textbf{What the module does.} This is the platform monitoring the company's own infrastructure --- an internal CMDB
plus an uptime/expiry watcher. It registers each \textbf{server} (a VPS or bare-metal host, with free-text capacity
fields for cpu/ram/storage/bandwidth and an \texttt{ip}) and each \textbf{service} hosted on a server (the API, PostgreSQL,
Redis, MinIO, with SSL/backup posture and a \texttt{domain}), records who is responsible via a \texttt{UserServerManagement}
join, and runs six per-minute cron jobs that do the actual monitoring. Three detector crons ICMP-ping
every running server, HTTP-probe every running service, and watch each row's \texttt{expiredAt} to warn before a paid
subscription lapses; three more are the delivery side. The design that ties detection to delivery is an
\textbf{outbox}: a detector never sends anything itself --- on a failed ping it inserts a \texttt{ServerNotification\{isSent:false\}}
row, and a \emph{separate} sender cron picks up the unsent rows, fans them out to each responsible manager's enabled
channels (reusing the exact \texttt{UserNotificationSettings} from Module A), and flips \texttt{isSent:true}. That
\texttt{isSent=false} flag doubles as natural dedup: the detector query skips any server that already has an undelivered
alert, so a persistently-down host re-alerts roughly once per delivery cycle rather than every minute. Every cron
runs under the same Postgres distributed lock the rest of the platform uses, so multiple API replicas can run
without double-alerting. The health-check cycle is Figure~\ref{fig:3-37}; the outbox-to-channels fan-out is
Figure~\ref{fig:3-38}.

Two facts about the module are worth stating plainly because they bound what it actually delivers. First,
\texttt{status} (Running/Stopped/Maintenance) is a \textbf{human-set label, not a live health state} --- the health checks
never write it; there is no \texttt{lastHealthCheck} column and no uptime history anywhere in the schema, so a check
leaves no persisted trace beyond a down-notification. Second, authorization is decided in the service,
not the guard: two role helpers gate every path --- the CTO gets global write and the CTO or CEO get global read,
while everyone else (in practice the DevOps engineer) is scoped by \texttt{server.managers} to the servers they manage.

% Figure 3.37 — render diagrams-src/fig_3_37.puml -> figures/fig_3_37.(png|pdf)
\reportfigure{figures/fig_3_37}{Server health-check cycle with outbox pattern.}{3-37}

% Figure 3.38 — render diagrams-src/fig_3_38.puml -> figures/fig_3_38.(png|pdf)
\reportfigure{figures/fig_3_38}{Alert fan-out from outbox to manager channels.}{3-38}

The class slice (Figure~\ref{fig:3-39}) is the module's five tables plus the shared status enum. \texttt{Server} and \texttt{Service}
mirror each other on billing/status/expiry fields; a \texttt{Service} cascade-deletes with its \texttt{Server}. The two
\texttt{*Notification} tables are the outbox rows. \texttt{UserServerManagement} is the responsibility join, and its scope is
worth noting: managers attach to \textbf{servers only} --- a service has no managers of its own, and its responsible
people are inherited from its parent server, a fact the fan-out relies on.

% Figure 3.39 — render diagrams-src/fig_3_39.puml -> figures/fig_3_39.(png|pdf)
\reportfigure{figures/fig_3_39}{Infrastructure monitoring class slice.}{3-39}

\textbf{Limitations.} The outbox pattern, lock-guarded crons and multi-channel reuse make this a well-structured
module, but two verified bugs undercut its core promise. \textbf{Service edit and delete are broken for the very role
meant to use them:} the non-CTO scoping builds a Prisma \texttt{where} of \texttt{\{ id, managers: \{ some: … \} \}} against the
\texttt{Service} model, which has no \texttt{managers} relation, so the query throws a validation error that escapes the
NotFound mapping and surfaces as a 500 --- and the DevOps engineer, who holds \texttt{service.update}/\texttt{service.delete}
and is not the CTO, hits it on every attempt. \textbf{The service health check almost certainly reports
every service as down:} \texttt{checkHttp} calls \texttt{axios.get} on a \emph{schemeless} domain (the seed stores
\texttt{api.tawer.tn}, not \texttt{https://api.tawer.tn}), which throws for lack of a protocol, is caught, and is treated as
a down result --- so a down-notification is created every cycle for healthy services. Around those,
the module has no alert cooldown or acknowledgement (a persistently-down or expiring resource re-alerts each
cycle), the \texttt{isSent} flip is fire-and-forget so a row is marked sent even if delivery failed, the four crons
carry \textasciitilde120 lines of copy-pasted band and fan-out logic that have already drifted (services lack the server's
final expiry band), and the specs are stubs --- one controller spec even imports a non-existent \texttt{UsersController}. There is also a frontend/backend RBAC divergence: the frontend grants CEO and customer-support
infra controls the API rejects, so those users see phantom buttons.

\subsection{Implementation}

The two modules present very differently in the UI. Notifications is woven through the whole shell rather than
sitting on a page: a bell dropdown in the header shows the unseen count and marks everything seen on open, a
full list page renders the inbox, and a dedicated settings page toggles email/Telegram/ntfy (push is registered
silently on first load, not toggled here) with the Telegram \texttt{chatId} input and the ntfy setup steps.
Push registration is automatic --- a header hook requests browser permission, obtains an FCM token, and registers
the device once per session. Infrastructure monitoring is two conventional list pages under an
\texttt{/infrastructure} route --- servers and services --- each a permission-gated table with an add/edit dialog and a
status badge; notably, there is \textbf{no live status or uptime indicator} in the UI, because the module persists no
health state (the badge shows the manual \texttt{status} label, and alerts arrive out-of-band over the channels, not on
the page).

The screens are supplied by the author:

\reportscreenshot{figures/screenshot_P27}{Notification bell dropdown and inbox list.}{P27}

\reportscreenshot{figures/screenshot_P28}{Notification channel settings page.}{P28}

\reportscreenshot{figures/screenshot_P29}{Infrastructure servers dashboard with add/edit dialog.}{P29}

\reportscreenshot{figures/screenshot_P30}{Services list with status badges and edit dialog.}{P30}

\subsection{Validation and testing}

The acceptance scenarios record the Given/When/Then behind each story. The genuine automated evidence in scope is \texttt{mail.service.spec.ts}, a unit test that exercises
the email channel. We scope the rest honestly --- coverage is otherwise near-zero on the backend: both
the notifications and the servers specs are \texttt{toBeDefined()} stubs (and the servers controller spec is outright
broken, importing a symbol that does not exist), so nothing tests the broadcast fan-out, the ownership scoping,
the six crons, the ping/HTTP probing, or the expiry-band math. The table below is therefore
verified by code reading, and we mark each qualification where the delivered behaviour diverges from the
story's intent --- two of the qualifications (ntfy dead, service checks false-positive) are the sprint's headline
bugs.

\begin{longtable}{@{}p{1.8cm}p{8.3cm}p{3.8cm}@{}}
\caption{Sprint 5 acceptance scenarios.}\label{tab:s5-acceptance} \\
\tableheaderrow
\thcell{US-ID} & \thcell{Acceptance scenario (Given / When / Then)} & \thcell{Result} \\
\midrule
\endfirsthead
\caption[]{Sprint 5 acceptance scenarios (continued).} \\
\tableheaderrow
\thcell{US-ID} & \thcell{Acceptance scenario (Given / When / Then)} & \thcell{Result} \\
\midrule
\endhead
US-S5-01 & Given a system event for a user, When it is created, Then a \texttt{UserNotification} row appears in that user's inbox, the bell shows an unseen count, and opening the list marks it seen (ownership-scoped). & \warn{} Verified by reading --- inbox persistence is coupled to \texttt{pushNotificationsEnabled}, so disabling push also loses in-app history \\
US-S5-02 & Given a user with a registered device, When a notification is sent, Then FCM pushes it to their tokens best-effort, surviving dead tokens. & \pass{} Verified by reading --- FCM multicast never throws; no test \\
US-S5-03 & Given a user, When they toggle email/push/Telegram/ntfy, Then subsequent system notifications respect those channel flags. & \warn{} Verified by reading --- email/push/Telegram honoured; \textbf{ntfy silently delivers nothing} because the topic is never written \\
US-S5-04 & Given a user, When they enter their Telegram \texttt{chatId}, Then alerts route to that chat via the shared bot token. & \pass{} Verified by reading --- per-user \texttt{chatId}, single shared bot secret; no test \\
US-S5-05 & Given a DevOps engineer or CTO, When they register a server (with managers) or a service, Then it is inventoried, manager-scoped for later access. & \warn{} Verified by reading --- create has no object-level authz and does not auto-add the creator as a manager, so the creator can lose write access to their own server \\
US-S5-06 & Given running servers/services, When the per-minute crons run, Then unreachable ones (ICMP/HTTP) produce a down-notification, guarded by a distributed lock. & \warn{} Verified by reading --- server ping path correct; \textbf{service HTTP probe uses a schemeless domain and reports every service down every cycle} \\
US-S5-07 & Given a down or expiring resource, When the sender cron runs, Then each responsible manager is alerted over their enabled channels and the outbox row is marked sent. & \warn{} Verified by reading --- delivery works, but the \texttt{isSent} flip is fire-and-forget (marked sent even on channel failure) and there is no cooldown/acknowledgement, so alerts can spam \\
US-S5-08 & Given the API is running, When \texttt{GET /health} is called without auth, Then it returns \texttt{\{status:'ok', timestamp\}}. & \pass{} Verified by reading --- single public liveness route, no guard \\
\bottomrule
\end{longtable}

\subsection{Sprint review}

This sprint made the platform's communication layer explicit and gave it its most demanding consumer. The
notifications module is now the documented substrate --- an in-app inbox, FCM push, per-user channel settings, and
Telegram/ntfy links --- that every other module reaches through one shared method for in-app/push and direct
channel calls for the rest. Infrastructure monitoring is the payoff: a self-observing loop where a cron detecting
a dead ping produces an alert that travels the whole backbone to a manager's Telegram, all reusing the same
settings. The outbox pattern, the distributed lock on every cron, and the correct ownership scoping on the
notification read/write surface are the sprint's real strengths.

Two of the sprint's features are provisioned but not yet wired end-to-end, and are prioritized first for the
hardening phase: ntfy delivery (the settings UI walks a user through subscribing to a topic that no code path
stores, so \texttt{NtfyService} early-returns and delivers nothing) and the service health-check (a schemeless-\texttt{domain}
probe that throws for the missing protocol and reads healthy services as ``down'', which without a cooldown would
flood managers). Those, together with the sprint's other notes --- restricting the
broadcast grant to executive/manager roles with throttling, correcting the non-CTO service scope through
\texttt{server.managers}, a single \texttt{NotificationDispatcherService} to replace the duplicated fan-out, decoupling the
in-app inbox from the push flag, a persisted \texttt{lastHealthCheckAt}, a unified channel error contract, and backend
tests for fan-out and RBAC scoping --- are consolidated with their remediations in the \emph{Limitations \& Perspectives}
discussion and the hardening backlog (Annex A).

\subsection{Cumulative class diagram}

Figure~\ref{fig:3-40} is the system after five sprints. Legibility now clearly outranks completeness: the five earlier
packages --- identity (Sprint 1), the project family (Sprint 2), the agile and task engines (Sprint 3), and the
Productivity Suite (Sprint 4) --- are shown plain and compacted to one or two identifying classes each, while
this sprint's two modules form a new \textbf{Communication \& Operations} package in yellow, split into its
notifications and monitoring halves. The seam that makes the picture cohere is the channel configuration:
\texttt{UserNotificationSettings} (with \texttt{UserTelegramBot} and \texttt{UserNtfyIntegration}) is the single place every
consumer reads before sending, so both the Sprint 4 reminder engine and this sprint's infrastructure alerts
depend on it. \texttt{UserServerManagement} joins \texttt{User} to \texttt{Server}, and the two outbox \texttt{*Notification} tables are
what the monitoring crons write and the fan-out reads.

% Figure 3.40 — render diagrams-src/fig_3_40.puml -> figures/fig_3_40.(png|pdf)
\reportfigure{figures/fig_3_40}{Cumulative class diagram after Sprint 5.}{3-40}
